\chapter{時間演化、守恆與耦合自旋}

量子態不是希爾伯特空間中的固定點。隨著時間推移，它會演化，但演化過程必須保留歸一化、機率，以及內積所表徵的可區分性。這些要求幾乎決定了量子動力學的形式。線性性帶來時間演化算符，機率守恆要求該算符為么正算符，而無窮小么正變換則由厄米算符生成。這個算符就是哈密頓量。我們將由這些陳述推導Schr\"odinger方程，探討能量本徵態如何累積相位，接著進行兩個具體計算：自旋在磁場中的進動，以及一對耦合自旋的兩個態之間的振盪。

\section{希爾伯特空間中的線性演化}

設時刻\(t\)的態為\(\ket{\psi(t)}\)。我們假設它與零時刻的態之間呈線性關係：
\begin{equation}
\ket{\psi(t)}=U(t)\ket{\psi(0)}.
\end{equation}
線性是一項公設，不是從更基本原理推導出的定理。它是量子力學中經實驗驗證的一部分；試圖取代它，通常會危及機率守恆或其他守恆定律。在選定的基底中，\(U(t)\)就是一個隨時間變化的矩陣。

對偶向量遵循共軛的演化規則。若
\begin{equation}
\ket{\phi(t)}=U(t)\ket{\phi(0)},
\end{equation}
則
\begin{equation}
\bra{\phi(t)}=\bra{\phi(0)}U^\dagger(t).
\end{equation}
對矩陣而言，匕首符號表示先取複共軛，再進行轉置。因此，左矢在對偶空間中遵循與右矢相同的演化規則。

\subsection{機率守恆蘊含么正性}

歸一化的態表示總機率為一：
\begin{equation}
\braket{\psi(t)}{\psi(t)}=1.
\end{equation}
時間演化必須保持這個範數。事實上，若線性映射保持每個向量的範數，就也會保持每個內積。從幾何上看，保持兩個向量及其差向量的長度，就會保留它們所構成的三角形，因而也保留兩者的夾角。在複向量空間中，再考慮帶有因子\(i\)的線性組合，也能得到同樣的結論。

因此，對任意的\(\ket{\phi}\)和\(\ket{\psi}\)，
\begin{equation}
\braket{\phi(t)}{\psi(t)}
=
\braket{\phi(0)}{\psi(0)}.
\end{equation}
代入演化關係可得
\begin{equation}
\bra{\phi(0)}U^\dagger(t)U(t)\ket{\psi(0)}
=
\braket{\phi(0)}{\psi(0)}.
\end{equation}
由於此式對所有向量對都成立，
\begin{equation}
\boxed{U^\dagger(t)U(t)=I}.
\end{equation}
這樣的算符稱為么正算符。其純量類比是模長為一的複數，
\begin{equation}
e^{i\theta}\bigl(e^{i\theta}\bigr)^*=1.
\end{equation}
么正算符是純相位的矩陣推廣：它可以使態在希爾伯特空間中旋轉，但不會將態拉伸或壓縮。

\section{無窮小生成元}

有限時間的演化可以由許多極小的時間步累積而成。設經過時間為\(\epsilon\)，小到可以忽略\(\epsilon^2\)階的項。經過時間為零時，演化算符就是單位算符：
\begin{equation}
U(0)=I.
\end{equation}
最低階的修正必須正比於\(\epsilon\)。按照慣例，我們將它寫成
\begin{equation}
\boxed{U(\epsilon)=I-\frac{i\epsilon}{\hbar}H}.
\end{equation}
在這個階段，只有因子\(\epsilon\)是必須如此選取的。\(i\)、\(\hbar\)以及\(H\)這個名稱，則是為了讓最後得到的算符具有標準的歸一化，並且是厄米算符。

\begin{classroomqa}
\lecturetimestamp{00:12:46}
\audiencequestion{為什麼用熟悉的單位表示時，普朗克常數小得如此驚人？}
\lecturerresponse{這個數值大小反映了人類選用的單位。化學家選擇了適合裝在手邊容器中的宏觀量，而原子比人小得多。在配合這些宏觀尺度的單位中，原子的量子作用量自然顯得極小。這不是另一個獨立的動力學謎題；很大程度上，這只是單位中所內含的人類歷史尺度。}
\end{classroomqa}

取厄米共軛，
\begin{equation}
U^\dagger(\epsilon)
=I+\frac{i\epsilon}{\hbar}H^\dagger.
\end{equation}
現在將兩者相乘，只保留一階項：
\begin{align}
U^\dagger(\epsilon)U(\epsilon)
&=
\left(I+\frac{i\epsilon}{\hbar}H^\dagger\right)
\left(I-\frac{i\epsilon}{\hbar}H\right) \\
&=
I+\frac{i\epsilon}{\hbar}\left(H^\dagger-H\right)
+O(\epsilon^2).
\end{align}
么正性要求這個乘積等於\(I\)，因此
\begin{equation}
\boxed{H^\dagger=H}.
\end{equation}
時間演化的無窮小生成元是厄米算符。由於厄米算符代表可觀測量，每個動力學系統都有一個生成其時間演化的可觀測量，這就是哈密頓量。我們以\(H\)表示該算符，並以\(E\)表示其可能的測量值，也就是本徵值。

\begin{classroomqa}
\lecturetimestamp{00:18:34}
\audiencequestion{為什麼這個生成元應該稱為能量？能量不就是\(p^2/2m\)，也就是做功的能力嗎？}
\lecturerresponse{這些說法在某些特殊的經典情境中很有用，但不是一般性的定義。表示式\(p^2/2m\)只是非相對論自由粒子的動能，無法用來定義例如儲存在電磁場中的能量。對抽象系統而言，首先要求能量具備的普遍特徵是守恆。我們很快就會證明\(H\)的期望值守恆，而其本徵態的相位也會滿足\(E=\hbar\omega\)。這兩項事實指出，這個生成元就是能量。}
\end{classroomqa}

\section{從微小步進到Schr\"odinger方程}

無窮小更新式其實已經包含了一個微分方程。對任意時間原點，
\begin{equation}
\frac{d}{dt}\ket{\psi(t)}
=
\lim_{\epsilon\to0}
\frac{\ket{\psi(t+\epsilon)}-\ket{\psi(t)}}{\epsilon}.
\end{equation}
黑板上先在\(t=0\)寫出這個式子；平移時間原點後，就得到上面的通式。

\begin{figure}[htbp]
\centering
\includegraphics[width=0.72\textwidth]{lecture_09_figure_03.png}
\caption{將無窮小演化轉化為微分方程的差商。}
\lectureframe{00:22:15}
\end{figure}

利用
\begin{equation}
\ket{\psi(t+\epsilon)}
=
\left(I-\frac{i\epsilon}{\hbar}H\right)\ket{\psi(t)}.
\end{equation}
可得
\begin{align}
\frac{d}{dt}\ket{\psi(t)}
&=
\lim_{\epsilon\to0}
\frac{\left(I-\frac{i\epsilon}{\hbar}H\right)\ket{\psi(t)}
-\ket{\psi(t)}}{\epsilon} \\
&=
-\frac{i}{\hbar}H\ket{\psi(t)}.
\end{align}
等價地，
\begin{equation}
\boxed{
i\hbar\frac{d}{dt}\ket{\psi(t)}
=H\ket{\psi(t)}
}.
\end{equation}
這就是Schr\"odinger方程。給定\(H\)及某一時刻的態，就能決定之後各時刻的態。之後進行測量時的機率，則可由一般的測量規則求得。

\section{能量本徵態與旋轉相位}

最簡單的解從哈密頓量的一個本徵向量開始：
\begin{equation}
H\ket{\psi_E}=E\ket{\psi_E}.
\end{equation}

\begin{figure}[htbp]
\centering
\includegraphics[width=0.70\textwidth]{lecture_09_figure_04.png}
\caption{黑板上的能量本徵值方程。}
\lectureframe{00:28:22}
\end{figure}

由於\(H\)作用在\(\ket{\psi_E}\)上仍得到同一個一維方向的向量，時間演化便不會使它與其他本徵向量混合；它只能乘上一個純量：
\begin{equation}
\ket{\psi_E(t)}=f(t)\ket{\psi_E}.
\end{equation}
代入Schr\"odinger方程可得
\begin{align}
i\hbar\,\dot f(t)\ket{\psi_E}
&=
f(t)H\ket{\psi_E} \\
&=
E f(t)\ket{\psi_E}.
\end{align}
因此
\begin{equation}
\dot f(t)=-\frac{iE}{\hbar}f(t),
\qquad
f(t)=e^{-iEt/\hbar},
\end{equation}
並且
\begin{equation}
\boxed{
\ket{\psi_E(t)}
=e^{-iEt/\hbar}\ket{\psi_E(0)}
}.
\end{equation}
能量本徵態之所以稱為定態，是因為它在希爾伯特空間中的射線不會移動；不過，它的右矢仍會累積相位。

指數函數可以寫成
\begin{equation}
e^{-iEt/\hbar}
=
\cos\!\left(\frac{Et}{\hbar}\right)
-i\sin\!\left(\frac{Et}{\hbar}\right).
\end{equation}
黑板推導時修正了負號，也補上了正弦函數內缺少的因子\(i\)。角頻率為
\begin{equation}
\omega=\frac{E}{\hbar},
\qquad
\boxed{E=\hbar\omega}.
\end{equation}
等價地，若使用一般頻率\(\nu=\omega/(2\pi)\)，則有\(E=h\nu\)。

\begin{figure}[htbp]
\centering
\includegraphics[width=0.78\textwidth]{lecture_09_figure_05.png}
\caption{相位解與能量--頻率關係分別在相鄰的黑板上推導。}
\lectureframe{00:36:19}
\end{figure}

\begin{classroomqa}
\lecturetimestamp{00:34:43}
\audiencequestion{為什麼在計算要求這麼做之前，就先在無窮小算符中加入因子\(i\)和\(\hbar\)？}
\lecturerresponse{這是預先考量答案所採用的慣例。若沒有\(i\)，么正性會要求生成元為反厄米算符。乘上\(i\)後，便得到本徵值為實數的厄米可觀測量。因子\(\hbar\)則使這些本徵值具有能量的單位及標準化慣例。物理上的要求是么正性；加入這些因子，是為了將其生成元寫成標準形式。}
\end{classroomqa}

\begin{classroomqa}
\lecturetimestamp{00:36:18}
\audiencequestion{\(E=\hbar\omega\)只適用於電磁波和光子嗎？}
\lecturerresponse{不是。光子是其中一個重要例子，但這個關係更普遍。任何哈密頓量本徵態都會獲得角頻率為\(E/\hbar\)的相位。要將電磁場模式的頻率與其光子的頻率連結起來，還需要更多物理內容；但能量--相位關係本身，已可由量子時間演化推出。}
\end{classroomqa}

\section{能量基底中的疊加}

單一能量本徵態只會改變整體相位，而整體相位不會影響機率。當一個態包含多個能量本徵向量時，有趣的動力學才開始。由於厄米哈密頓量具有完備的正交歸一本徵基底，
\begin{equation}
H\ket{\psi_n}=E_n\ket{\psi_n},
\qquad
\ket{\psi(0)}=\sum_n\alpha_n\ket{\psi_n}.
\end{equation}

\begin{figure}[htbp]
\centering
\includegraphics[width=0.72\textwidth]{lecture_09_figure_06.png}
\caption{任意態在哈密頓量本徵向量上的展開。}
\lectureframe{00:38:35}
\end{figure}

由線性性可知，各分量可以分別演化：
\begin{equation}
\boxed{
\ket{\psi(t)}
=
\sum_n\alpha_n e^{-iE_n t/\hbar}\ket{\psi_n}
}.
\end{equation}
高能分量在複平面上快速旋轉，低能分量則轉得較慢。會改變可測機率的，不是整個態共同具有的相位，而是能量\(E_n\)不同的分量之間不斷變化的相對相位。

在能量基底中，
\begin{equation}
\begin{aligned}
H&=
\begin{pmatrix}
E_1&0&\cdots\\
0&E_2&\cdots\\
\vdots&\vdots&\ddots
\end{pmatrix},\\[0.6em]
U(t)&=
\begin{pmatrix}
e^{-iE_1t/\hbar}&0&\cdots\\
0&e^{-iE_2t/\hbar}&\cdots\\
\vdots&\vdots&\ddots
\end{pmatrix}.
\end{aligned}
\end{equation}
若哈密頓量不隨時間變化，也可以不選定基底，直接概括為
\begin{equation}
\boxed{U(t)=e^{-iHt/\hbar}},
\end{equation}

\begin{classroomqa}
\lecturetimestamp{00:40:17}
\audiencequestion{這種分解本質上是傅立葉級數嗎？它是否類似於熱傳導問題中使用的分解？}
\lecturerresponse{確實可以說，這是在時間上的類傅立葉分解：每個能量本徵向量都帶有確定的頻率。傅立葉與拉普拉斯方法也能將許多經典方程中的模式分離，包括擴散方程。兩者在數學上的相似性確實存在，但頻率對應哪一種物理量，必須逐一檢視。在量子力學中，具體關係由Schr\"odinger方程決定。}
\end{classroomqa}

\subsection{系統組合}

對兩個互不作用的系統\(A\)和\(B\)，組合系統的哈密頓量是在乘積空間上作用的兩個生成元之和：
\begin{equation}
H_{AB}
=
H_A\otimes I_B+I_A\otimes H_B.
\end{equation}
若
\begin{equation}
H_A\ket{a}=E_a\ket{a},
\qquad
H_B\ket{b}=E_b\ket{b},
\end{equation}
則
\begin{equation}
H_{AB}\ket{a,b}=(E_a+E_b)\ket{a,b}.
\end{equation}
對彼此分離且互不作用的系統而言，能量會相加，因此相位頻率也會相加。相互作用會額外貢獻一項\(H_{\mathrm{int}}\)，此時這個簡單的和便不再代表完整的能量。

\begin{classroomqa}
\lecturetimestamp{00:56:20}
\audiencequestion{如果兩個分離的系統各自有頻率相同的分量，將兩個系統組合起來時會發生什麼？}
\lecturerresponse{對互不作用的系統而言，組合系統的能量本徵值是各系統本徵值的和。頻率相同可能造成簡併，但不會因此使相應的本徵向量混合。要產生混合，必須有相互作用項，或是在簡併子空間中選擇另一組基底。系統開始互相作用後，就必須把擴大的哈密頓量視為新問題重新對角化。}
\end{classroomqa}

\section{期望值與對易子}

我們在保持可觀測量不變的同時，追蹤態的演化。這就是Schr\"odinger繪景。設\(A\)為一個沒有顯式時間依賴的可觀測量。它的期望值為
\begin{equation}
\langle A\rangle_t
=
\bra{\psi(t)}A\ket{\psi(t)}.
\end{equation}
分別對左矢和右矢微分：
\begin{equation}
\frac{d}{dt}\langle A\rangle_t
=
\bra{\dot\psi}A\ket{\psi}
+
\bra{\psi}A\ket{\dot\psi}.
\end{equation}
Schr\"odinger方程及其厄米共軛給出
\begin{equation}
\ket{\dot\psi}
=-\frac{i}{\hbar}H\ket{\psi},
\qquad
\bra{\dot\psi}
=\frac{i}{\hbar}\bra{\psi}H.
\end{equation}
因此
\begin{align}
\frac{d}{dt}\langle A\rangle_t
&=
\frac{i}{\hbar}\bra{\psi}HA\ket{\psi}
-\frac{i}{\hbar}\bra{\psi}AH\ket{\psi} \\
&=
\frac{i}{\hbar}\bra{\psi}(HA-AH)\ket{\psi}.
\end{align}
定義對易子
\begin{equation}
[H,A]=HA-AH.
\end{equation}
期望值的運動方程便是
\begin{equation}
\boxed{
\frac{d}{dt}\langle A\rangle_t
=
\frac{i}{\hbar}\langle[H,A]\rangle_t
}.
\end{equation}

\begin{figure}[htbp]
\centering
\includegraphics[width=0.74\textwidth]{lecture_09_figure_07.png}
\caption{期望值方程，以及對易子\(AB-BA=[A,B]\)。}
\lectureframe{00:51:05}
\end{figure}

\begin{classroomqa}
\lecturetimestamp{00:49:13}
\audiencequestion{為什麼這個方程中的對易子也要放在期望值裡？}
\lecturerresponse{對易子是一個算符。任何放在\(\bra{\psi}\)與\(\ket{\psi}\)之間的算符，都有相應的期望值；因此\(\bra{\psi}[H,A]\ket{\psi}\)就是該對易子在目前態中的平均值。}
\end{classroomqa}

\subsection{守恆定律}

令\(A=H\)。算符總是與自身對易：
\begin{equation}
[H,H]=HH-HH=0.
\end{equation}
因此
\begin{equation}
\boxed{
\frac{d}{dt}\langle H\rangle_t=0
}.
\end{equation}
哈密頓量的期望值守恆。這就是我們預告過、將\(H\)認定為能量的結構性理由。每個孤立系統都有哈密頓量，而每個不隨時間變化的哈密頓量都帶來這項守恆定律。

更一般而言，
\begin{equation}
[H,A]=0
\quad\Longrightarrow\quad
\frac{d}{dt}\langle A\rangle_t=0.
\end{equation}

\begin{classroomqa}
\lecturetimestamp{00:51:41}
\audiencequestion{與\(H\)對易的所有算符，代表哪些物理量？其中是否也有能量以外的量？}
\lecturerresponse{這些算符代表系統的守恆可觀測量。有哪些守恆量，取決於哈密頓量。自由粒子的動量守恆，但力的作用可能使動量不再守恆。能量必定是守恆量，因為\(H\)與自身對易；其他與哈密頓量對易的算符，則代表額外的對稱性與守恆律。}
\end{classroomqa}

\section{磁場中的單一自旋}

考慮一個自旋，其泡利算符為
\begin{equation}
\boldsymbol{\sigma}
=(\sigma_1,\sigma_2,\sigma_3)
\end{equation}
並處於固定的外加磁場\(\mathbf B\)中。磁矩與自旋成正比。依照本文採用的慣例，
\begin{equation}
\boxed{
H=\frac{\mu}{2}\boldsymbol{\sigma}\cdot\mathbf B
}.
\end{equation}
因子\(1/2\)、正負號，以及依電荷而定的尺度，都可以吸收到\(\mu\)的定義中。磁場視為經典場：由一個大型且穩定的磁鐵提供\(\mathbf B\)，我們不會對該磁鐵的量子態建模。

\begin{figure}[htbp]
\centering
\includegraphics[width=0.68\textwidth]{lecture_09_figure_08.png}
\caption{磁矩算符與外加磁場耦合。}
\lectureframe{01:01:45}
\end{figure}

令磁場沿第三軸，
\begin{equation}
\mathbf B=B\hat{\mathbf z},
\qquad
H=\frac{\mu B}{2}\sigma_3.
\end{equation}
我們將追蹤自旋三個分量的期望值。為了簡化方程，記為
\begin{equation}
s_i(t)=\langle\sigma_i\rangle_t.
\end{equation}
縱向分量滿足
\begin{equation}
\dot s_3
=
\frac{i}{\hbar}\langle[H,\sigma_3]\rangle
=0,
\end{equation}
這是因為\(H\)正比於\(\sigma_3\)。自旋沿磁場方向的投影保持不變。從幾何上看，平均自旋的端點必須留在以磁場軸為軸、頂角固定的圓錐面上。

\begin{figure}[htbp]
\centering
\includegraphics[width=0.58\textwidth]{lecture_09_figure_09.png}
\caption{縱向自旋保持不變，因此運動沿著以磁場軸為軸的圓錐面進行。}
\lectureframe{01:04:45}
\end{figure}

對第一個橫向分量，
\begin{equation}
\dot s_1
=
\frac{i\mu B}{2\hbar}
\langle[\sigma_3,\sigma_1]\rangle.
\end{equation}
可以直接利用下式驗算這個對易子：
\begin{equation}
\begin{aligned}
\sigma_3&=
\begin{pmatrix}1&0\\0&-1\end{pmatrix},
&
\sigma_1&=
\begin{pmatrix}0&1\\1&0\end{pmatrix},\\[0.5em]
\sigma_2&=
\begin{pmatrix}0&-i\\i&0\end{pmatrix}.
\end{aligned}
\end{equation}
確實，
\begin{equation}
\sigma_3\sigma_1=i\sigma_2,
\qquad
\sigma_1\sigma_3=-i\sigma_2,
\end{equation}
所以
\begin{equation}
[\sigma_3,\sigma_1]=2i\sigma_2.
\end{equation}
因此，
\begin{equation}
\dot s_1=-\frac{\mu B}{\hbar}s_2.
\end{equation}
同理，利用\([\sigma_3,\sigma_2]=-2i\sigma_1\)，
\begin{equation}
\dot s_2=\frac{\mu B}{\hbar}s_1.
\end{equation}

\begin{figure}[htbp]
\centering
\includegraphics[width=0.72\textwidth]{lecture_09_figure_10.png}
\caption{推導橫向運動方程時的泡利矩陣對易子計算。}
\lectureframe{01:09:35}
\end{figure}

令
\begin{equation}
\Omega=\frac{\mu B}{\hbar}.
\end{equation}
三個方程為
\begin{equation}
\dot s_1=-\Omega s_2,
\qquad
\dot s_2=\Omega s_1,
\qquad
\dot s_3=0.
\end{equation}
橫向振幅為\(s_\perp\)的一個解是
\begin{equation}
s_1(t)=s_\perp\cos(\Omega t+\varphi),
\qquad
s_2(t)=s_\perp\sin(\Omega t+\varphi),
\end{equation}
其中\(s_3\)為常數。

\begin{figure}[htbp]
\centering
\includegraphics[width=0.64\textwidth]{lecture_09_figure_11.png}
\caption{橫向自旋分量形成一對正弦--餘弦函數，其角頻率為\(\mu B/\hbar\)。}
\lectureframe{01:11:45}
\end{figure}

因此，平均自旋會像陀螺一樣繞著磁場進動。選擇磁場沿\(z\)方向只是為了簡化計算。若磁場方向任意，平行於\(\mathbf B\)的分量仍守恆，而垂直分量則繞著它旋轉。

\begin{classroomqa}
\lecturetimestamp{01:13:02}
\audiencequestion{正弦與餘弦解是否應該多帶一個相位？}
\lecturerresponse{是的。相位\(\varphi\)指定橫向平面中的初始方向；等價地說，它指定哪一個時刻被定義為\(t=0\)。將相位設為零，只是選定自旋從某個特定的橫向軸方向開始。}
\end{classroomqa}

\subsection{態的製備與測量}

假設一個系綜被製備在\(\sigma_1\)的\(+1\)本徵態中。在\(t=0\)，測量\(\sigma_1\)必定得到\(+1\)；測量\(\sigma_2\)或\(\sigma_3\)時，得到\(\pm1\)的兩種結果之機率則相等。經過四分之一個進動週期後，態已旋轉至\(\sigma_2\)的\(+1\)本徵態：此時測量\(\sigma_2\)必定得到確定結果，而\(\sigma_1\)的測量結果則完全不確定。再經過四分之一個週期，態便到達\(\sigma_1\)的\(-1\)本徵態。因此，旋轉的期望值具有直接的操作意義。

\begin{classroomqa}
\lecturetimestamp{01:14:04}
\audiencequestion{若期望值需要對許多電子進行統計，系綜中的電子不會各自帶有不同相位，因而使進動彼此抵銷嗎？}
\lecturerresponse{若對每個電子進行相同的實驗，就不會。相同的製備會使每個電子都有相同的初始態，因此相位演化也相同。若初始相位未受控制，這個系綜就不是純態製備，而是混合態；退相干可能會使橫向平均值減小。}
\end{classroomqa}

\section{兩個耦合自旋}

現在來看兩個自旋。光說它們彼此糾纏，並不足以預測它們將如何演化。要描述動力學，還需要哈密頓量。

\begin{classroomqa}
\lecturetimestamp{01:16:49}
\audiencequestion{如果兩個電子處於糾纏態，沿任意共同軸測量都會得到相反結果，這是否表示它們的運動總是相差\(180^\circ\)的相位？}
\lecturerresponse{單憑糾纏無法得出動力學上的結論。題述的反相關描述的是單重態，但之後如何演化取決於能量算符。我們必須先說明是否存在外加磁場、自旋間的相互作用，或兩者兼有。}
\end{classroomqa}

考慮一種簡單的等向性交互作用，其中每個自旋都提供另一個自旋所感受到的有效場：
\begin{equation}
\boxed{
H=J\,\boldsymbol{\sigma}\cdot\boldsymbol{\tau}
=
J\left(
\sigma_1\tau_1+\sigma_2\tau_2+\sigma_3\tau_3
\right)
}.
\end{equation}
算符\(\boldsymbol{\sigma}\)作用在第一個自旋上，\(\boldsymbol{\tau}\)則作用在第二個自旋上。耦合常數\(J\)設定能量尺度；此處不指定其數值與正負號。

\begin{figure}[htbp]
\centering
\includegraphics[width=0.75\textwidth]{lecture_09_figure_12.png}
\caption{等向性雙自旋交互作用\(H\propto\boldsymbol{\sigma}\cdot\boldsymbol{\tau}\)。}
\lectureframe{01:19:20}
\end{figure}

\begin{classroomqa}
\lecturetimestamp{01:17:49}
\audiencequestion{將兩個糾纏電子置於均勻磁場中，是否必然會使它們失去糾纏？}
\lecturerresponse{不一定。磁場會改變哈密頓量，也可能使自旋傾向於對齊，尤其是磁場很強時；然而，均勻磁場中的么正演化本身並不保證糾纏會消失。答案取決於完整的哈密頓量與初始態。}
\end{classroomqa}

\subsection{三重態與單重態的能量}

雙自旋系統的希爾伯特空間維度為四。此交互作用在三重態--單重態基底中是對角的：
\begin{align}
\ket{T_+}&=\ket{\uparrow\uparrow}, \\
\ket{T_0}&=
\frac{\ket{\uparrow\downarrow}+\ket{\downarrow\uparrow}}{\sqrt2}, \\
\ket{T_-}&=\ket{\downarrow\downarrow}, \\
\ket{S}&=
\frac{\ket{\uparrow\downarrow}-\ket{\downarrow\uparrow}}{\sqrt2}.
\end{align}
三個三重態具有相同的能量\(E_T\)；單重態則有另一個能量\(E_S\)。因此，雖有四個線性獨立的本徵向量，卻只有兩個不同的能階。三重態能階為三重簡併。

\begin{figure}[htbp]
\centering
\includegraphics[width=0.74\textwidth]{lecture_09_figure_13.png}
\caption{耦合自旋的哈密頓量有一個簡併的三重態能階，以及另一個單重態能階。}
\lectureframe{01:20:30}
\end{figure}

\begin{classroomqa}
\lecturetimestamp{01:21:16}
\audiencequestion{如果初始態本來就是單重態能量本徵態，會如何演化？}
\lecturerresponse{它會保持在單重態，只累積相位\(e^{-iE_St/\hbar}\)。如同任何單一能量本徵態，僅僅因為整個右矢乘上這個共同相位，其可測量的性質並不會改變。}
\end{classroomqa}

\begin{classroomqa}
\lecturetimestamp{01:21:54}
\audiencequestion{在這項計算中，外加磁場是否向雙自旋系統提供能量？}
\lecturerresponse{這裡求解的例子沒有外加磁場。此處唯一討論的場，是\(J\,\boldsymbol{\sigma}\cdot\boldsymbol{\tau}\)所表示的自旋間交互作用場。加入外加磁場會在\(H\)中增添另一項，因而構成不同的問題。}
\end{classroomqa}

\subsection{從一個自旋向上、另一個自旋向下開始}

能量本徵態只會累積相位，因此我們選擇一個不是能量本徵態的直積態：
\begin{equation}
\ket{\psi(0)}=\ket{\uparrow\downarrow}.
\end{equation}
將它展開為兩個相關的本徵向量：
\begin{equation}
\boxed{
\ket{\uparrow\downarrow}
=
\frac{1}{\sqrt2}\ket{T_0}
+
\frac{1}{\sqrt2}\ket{S}
}.
\end{equation}

\begin{figure}[htbp]
\centering
\includegraphics[width=0.76\textwidth]{lecture_09_figure_14.png}
\caption{\(\ket{\uparrow\downarrow}\)這個直積態分解為三重態與單重態分量。}
\lectureframe{01:24:30}
\end{figure}

現在，各能量分量分別累積自己的相位：
\begin{equation}
\ket{\psi(t)}
=
\frac{e^{-iE_Tt/\hbar}}{\sqrt2}\ket{T_0}
+
\frac{e^{-iE_St/\hbar}}{\sqrt2}\ket{S}.
\end{equation}

\begin{classroomqa}
\lecturetimestamp{01:25:59}
\audiencequestion{三重態能量\(E_T\)和單重態能量\(E_S\)，哪一個較大？}
\lecturerresponse{這取決於耦合常數的正負號。若\(H=J\,\boldsymbol{\sigma}\cdot\boldsymbol{\tau}\)且\(J>0\)，單重態能量較低，三重態能量較高；耦合常數反號時，能量排序便會顛倒。若對兩個能量都加上同一常數，只會改變整體相位，因此可以把任一能階設為零，而不影響躍遷機率。}
\end{classroomqa}

回到直積基底：
\begin{align}
\ket{\psi(t)}
&=
\frac12
\left(
e^{-iE_Tt/\hbar}+e^{-iE_St/\hbar}
\right)\ket{\uparrow\downarrow} \\
&\quad+
\frac12
\left(
e^{-iE_Tt/\hbar}-e^{-iE_St/\hbar}
\right)\ket{\downarrow\uparrow}.
\end{align}

\begin{figure}[htbp]
\centering
\includegraphics[width=0.76\textwidth]{lecture_09_figure_15.png}
\caption{演化後的態重新展開於\(\ket{\uparrow\downarrow}\)、\(\ket{\downarrow\uparrow}\)基底中。}
\lectureframe{01:30:40}
\end{figure}

在\(t=0\)時，\(\ket{\downarrow\uparrow}\)的係數為零，這是必然的。之後的時刻，單重態與三重態分量之間的相對相位會使這個係數增大。定義平均能量與能量差，
\begin{equation}
\overline E=\frac{E_T+E_S}{2},
\qquad
\Delta E=E_T-E_S.
\end{equation}
提出平均能量的相位因子後，可得
\begin{equation}
\boxed{
\begin{aligned}
\ket{\psi(t)}
&=e^{-i\overline E t/\hbar}\Biggl[
\cos\!\left(\frac{\Delta E\,t}{2\hbar}\right)
\ket{\uparrow\downarrow}
\\
&\qquad
-i\sin\!\left(\frac{\Delta E\,t}{2\hbar}\right)
\ket{\downarrow\uparrow}
\Biggr]
\end{aligned}
}.
\end{equation}

\begin{figure}[H]
\centering
\includegraphics[width=0.76\textwidth]{lecture_09_figure_16.png}
\caption{平均能量只帶來共同相位；單重態--三重態的能量差則出現在正弦與餘弦振幅中。}
\lectureframe{01:35:15}
\end{figure}

因子\(e^{-i\overline E t/\hbar}\)對整個右矢都是共同的，會在所有機率中消去。可觀測的機率為
\begin{equation}
P_{\uparrow\downarrow}(t)
=
\cos^2\!\left(\frac{\Delta E\,t}{2\hbar}\right),
\qquad
P_{\downarrow\uparrow}(t)
=
\sin^2\!\left(\frac{\Delta E\,t}{2\hbar}\right).
\end{equation}
這個態會在\(\ket{\uparrow\downarrow}\)和\(\ket{\downarrow\uparrow}\)之間來回振盪，振盪速率只由能量差決定。在中間時刻，兩個振幅都不為零；除少數特殊時刻外，兩個自旋都處於糾纏態。

這項計算展示了一種通用方法：先選定適合當前物理問題的基底，再將初始態展開為哈密頓量的本徵向量；讓各本徵向量分別累積相位，最後在測量基底中重新組合各分量。哈密頓量一旦對角化，量子動力學的處理便有章可循；所有可觀測的演化都由相對相位造成。
